% Fig — the assembled tridiagonal system: banded structure + the two
% boundary rows that differ. Body only.
\begin{tikzpicture}[font=\footnotesize, every node/.style={align=center},
  lbl/.style={font=\scriptsize\color{dim}}]
  \def\n{8}
  \foreach \i in {0,...,7} \foreach \j in {0,...,7} {
    \pgfmathtruncatemacro\d{abs(\i-\j)}
    \ifnum\d<2
      \pgfmathsetmacro\y{-0.62*\i} \pgfmathsetmacro\x{0.62*\j}
      \ifnum\i=0 \fill[coral!75] (\x,\y) circle (2.6pt);
      \else\ifnum\i=7 \fill[forest!80] (\x,\y) circle (2.6pt);
      \else \fill[char!75] (\x,\y) circle (2.6pt); \fi\fi
    \fi }
  \draw[rounded corners=2pt, color=dim, line width=0.7pt] (-0.35,0.35) rectangle (4.7,-4.7);
  \node[lbl, above=6pt] at (2.2,0.3) {$A\,T^{\,n+1}=d$ \; (one solve per time step, $O(N)$ by Thomas)};
  % row annotations
  \node[font=\scriptsize\color{coral}, anchor=west, text width=58mm] at (5.3,0.0)
    {\textbf{row 0 --- Dirichlet ghost:} $T_s$ is already known from Newton,
     so its flux moves to the RHS:
     $d_0\to d_0-\tfrac12\alpha T_0^{\,n}
      +\tfrac12\alpha(T_s^{\,n}{+}T_s^{\,n+1})$};
  \node[font=\scriptsize\color{char}, anchor=west, text width=58mm] at (5.3,-2.2)
    {\textbf{interior rows:} the pure CN stencil ---
     $(a_i,\,b_i,\,c_i)$; in the uniform toy
     $(-\tfrac12\alpha,\;1{+}\alpha,\;-\tfrac12\alpha)$};
  \node[font=\scriptsize\color{forest}, anchor=west, text width=58mm] at (5.3,-4.35)
    {\textbf{row $n{-}1$ --- Neumann ghost:} the basal flux ---
     $b_{n-1}\to b_{n-1}-\tfrac12\alpha$,\;
     $d_{n-1}\to d_{n-1}+\Delta t\,Q_b/C$};
\end{tikzpicture}
